
\phantomsection
\addchaptertocentry{Удиртгал}
\label{Introduction}

\section*{Удиртгал}
Сүүлийн жилүүдэд цахим сургалт нь боловсролын байгууллага, хувийн сургалтын төв, бие даан суралцагчдын хувьд үндсэн хэрэгсэл болон хөгжиж байна. Интернэт холболт, веб технологи, үүлэн тооцоолол, мультимедиа агуулгын хүртээмж нэмэгдэхийн хэрээр сургалтын материалыг хүссэн цагтаа, хүссэн газраасаа ашиглах боломж улам өргөжив. Гэсэн хэдий ч цахим сургалтын олон платформ нийтлэг бүтэцтэй хэвээр байгаа бөгөөд хэрэглэгч бүрт ижил урсгал, ижил хичээлийн дараалал, ижил үнэлгээний механизм санал болгодог. Энэ нь суралцагчдын анхны мэдлэгийн түвшин, суралцах хэв маяг, хурдац, зорилго харилцан адилгүй байдаг бодит нөхцөлтэй бүрэн нийцдэггүй.

Ялангуяа програмчлалын сургалтын салбарт энэхүү асуудал илүү тод илэрдэг. Нэг хэрэглэгч хувьсагч, өгөгдлийн төрөл, нөхцөлт салбарлалт зэрэг суурь ойлголтыг дахин бататгах шаардлагатай байхад нөгөө хэрэглэгч функц, өгөгдлийн бүтэц, веб хөгжүүлэлтийн суурь ойлголт руу шууд шилжих боломжтой байдаг. Хэрэв систем эдгээр ялгааг тооцохгүй бол анхан түвшний хэрэглэгч агуулгыг хэт хурдан гэж үзэж, ахисан түвшний хэрэглэгч агуулгыг хэт удаан, нэгэн хэвийн гэж мэдрэх эрсдэлтэй. Үүний үр дүнд суралцах идэвх буурах, хичээлээ гүйцээхгүй орхих, системийг тогтмол ашиглахгүй болох зэрэг сөрөг үзэгдэл нэмэгддэг.

Нөгөө талаас хиймэл оюуны сүүлийн үеийн хөгжил, ялангуяа том хэлний загваруудын практик хэрэглээ нь цахим сургалтыг илүү ухаалаг, илүү хувь хүнд тохирсон хэлбэрт шилжүүлэх боломжийг нээж байна. Хичээлийн агуулгыг автоматаар хураангуйлах, ойлголт шалгах асуулт боловсруулах, суралцагчийн асуултад байгалийн хэлээр хариулах, нэмэлт практик даалгавар санал болгох зэрэг нь өмнө нь ихэвчлэн багшийн гар ажиллагаанд тулгуурлаж байсан үйлдлүүд юм. Харин одоо эдгээрийг хиймэл оюуны үйлчилгээтэй уялдуулан бодит системийн түвшинд хэрэгжүүлэх боломж бүрдээд байна.

Дасан зохицох сургалт буюу хэрэглэгчийн түвшин, ахиц, гүйцэтгэл, сонирхолд үндэслэн агуулгыг дасан зохицуулах арга нь энэ өөрчлөлтийн гол цөмд оршдог. Ийм систем нь зөвхөн ``ямар контент байна'' гэдэгт төвлөрөх бус ``хэн суралцаж байна'', ``ямар түвшинд байна'', ``дараа нь юу сурах ёстой вэ'', ``аль хэсэгт тусламж хэрэгтэй байна'' гэсэн асуултад хариулдаг байх шаардлагатай. Өөрөөр хэлбэл цахим сургалтын платформ нь статик контент түгээгч хэрэгсэл бус, суралцах үйл явцыг удирдах ухаалаг орчин болох ёстой.

Энэхүү дипломын ажлын хүрээнд програмчлалын сургалтад зориулсан, хиймэл оюунд суурилсан, шаталсан сургалтын веб системийг судалж, зохион бүтээж, хөгжүүлэв. Уг систем нь хэрэглэгчийн анхны мэдлэгийг түвшин тогтоох шалгалтаар тодорхойлох, түүнд тохирсон курс болон хичээл санал болгох, хичээлийн түвшинд хиймэл оюунаар үүсгэсэн хураангуй, сорил, дадлагын ажил үзүүлэх, ахицын хураангуйгаар суралцах явцыг тоон үзүүлэлтээр хянах, түвшин ахих шалгалтаар дараагийн түвшинд ахих боломж олгох зорилготой. Үүнээс гадна хэрэглэгчдэд харилцан санал солилцох хамтын орчин, сонирхолтой байдлыг нэмэгдүүлэх тоглоомын хэсэг, системийн түвшний чатбот туслах зэрэг боломжийг нэг дор нэгтгэсэн.

Ажлын онцлог нь зөвхөн интерфейсийн загвар гаргах эсвэл онолын түвшинд санаа дэвшүүлэхэд хязгаарлагдахгүй, харин урд талын хэсэг, арын хэсэг, өгөгдлийн сан, хиймэл оюуны уялдуулалт, нэвтрүүлэлтийн орчин, нөөц логик, хяналтын механизм зэрэг бодит системийн бүрэлдэхүүнүүдийг цогцоор нь хэрэгжүүлсэнд оршино. Энэ нь дипломын ажлыг зөвхөн судалгааны тайлан бус, бодит хэрэглээний боломжтой инженерчлэлийн шийдэл болгон хөгжүүлэх суурь болсон.

\subsection*{Сэдвийг сонгосон үндэслэл}
Програмчлалын боловсролын хэрэгцээ тогтмол өсч байгаа боловч бүх суралцагч ижил эхлэлтэй байдаггүй. Зарим нь анхан шатны ойлголтоо бататгах шаардлагатай байхад зарим нь тодорхой технологийн стек рүү шууд шилжих боломжтой байдаг. Ийм нөхцөлд хэрэглэгч бүрт ижил агуулга түгээх систем нь сургалтын үр өгөөжийг бууруулна. Мөн Монгол хэл дээр ойлгомжтой тайлбар, товч хураангуй, асуулт-хариултын дэмжлэг үзүүлэх ухаалаг системийн хэрэгцээ өндөр хэвээр байна.

Нөгөө талаар сургалтын агуулга бэлтгэх, асуулт боловсруулах, нэмэлт даалгавар гаргах, хэрэглэгчийн ахицыг хянах зэрэг ажил нь багш, агуулга бэлтгэгчдээс ихээхэн хугацаа шаарддаг. Хиймэл оюун ашигласнаар эдгээр үйл явцын тодорхой хэсгийг автоматжуулж, агуулгын хүртээмж, хэрэглэгчийн оролцоо, тайлбарын хурдыг сайжруулах боломжтой. Иймээс дасан зохицох сургалт болон үүсгэгч хиймэл оюуныг хослуулсан систем боловсруулах нь судалгааны болон практик ач холбогдолтой сэдэв юм.

\subsection*{Зорилго}
Програмчлалын сургалтад зориулсан, хэрэглэгчийн түвшин болон ахицад тулгуурлан агуулга санал болгодог, хиймэл оюуны дэмжлэгтэй шаталсан сургалтын веб системийг судалж, хөгжүүлэх.

\subsection*{Зорилтууд}
\begin{enumerate}[nosep]
\item Цахим сургалт, дасан зохицох сургалт, хувь хүнд тохирсон зөвлөмж, үүсгэгч хиймэл оюуны сургалтын хэрэглээний онолын үндсийг судлах
\item Төслийн хэрэглэгчийн хэрэгцээ, функциональ болон функцийн бус шаардлагыг тодорхойлох
\item Урд талын хэсэг, арын хэсэг, өгөгдлийн сан, хиймэл оюуны үйлчилгээ болон нэвтрүүлэлтийн орчныг багтаасан системийн ерөнхий архитектур боловсруулах
\item Хэрэглэгчийн түвшин тогтоох шалгалт болон түвшин ахиулах шалгалтын логикийг хэрэгжүүлэх
\item Курс, хичээл, суралцах замнал, ахиц хяналт, таалагдсан хичээлүүд болон зөвлөмжийн урсгалыг хөгжүүлэх
\item Хиймэл оюунаар үүсгэсэн хураангуй, сорил, дадлагын даалгавар, чатбот, сэтгэгдлийн хяналт зэрэг ухаалаг боломжуудыг системд нэгтгэх
\item Хөгжүүлсэн системийг туршиж, гарсан үр дүн, давуу тал, хязгаарлалтыг дүгнэх
\end{enumerate}

\subsection*{Судалгааны объект ба судлагдахуун}
Судалгааны объект нь вебэд суурилсан програмчлалын сургалтын платформ болно. Судлагдахуун нь тухайн платформ доторх шаталсан сургалтын логик, хэрэглэгчийн түвшин тогтоох ба ахиц хянах механизм, хиймэл оюуны тусламжтай сургалтын үйлдлүүд, хэрэглэгчийн туршлагыг сайжруулах урд ба арын хэсгийн уялдаа юм.

\subsection*{Ажлын практик ач холбогдол}
Энэхүү ажил нь дараах практик ач холбогдолтой.
\begin{itemize}
\item Програмчлалын сургалтын агуулгыг хэрэглэгчийн түвшинд илүү оновчтой түгээх системийн загвар бий болгоно.
\item Хиймэл оюун ашиглан хичээлийн тайлбар, шалгалтын асуулт, практик даалгаврыг дэмжих бодит аргачлал санал болгоно.
\item Урд талын хэсэг, арын хэсэг, өгөгдлийн сан, хиймэл оюуны уялдуулалт бүхий бүтэн стек систем хөгжүүлэх жишиг хэрэгжилт болно.
\item Цаашид багшийн самбар, илүү нарийвчилсан зөвлөмжийн систем, ахицын шинжилгээ, гар утсанд нийцсэн хувилбар зэрэг өргөтгөл хийх суурь платформ бүрдүүлнэ.
\end{itemize}

\subsection*{Шинэлэг тал}
Хөгжүүлсэн системийн шинэлэг талыг дараах байдлаар тодорхойлж болно.
\begin{enumerate}[nosep]
\item Түвшин тогтоох шалгалт болон түвшин ахих шалгалтыг нэгтгэсэн шаталсан сургалтын урсгалтай
\item Хичээлийн түвшинд хиймэл оюунаар үүсгэсэн хураангуй, сорил, дадлагын даалгавар ашигладаг
\item Хэрэглэгчийн явц, оноо, долоо хоногийн идэвхийг ахицын хураангуй хэсэгт нэгтгэн харуулдаг
\item Хичээлтэй холбоотой хэлэлцүүлгийг хамтын орчны модульд нэгтгэж, хяналтын логикоор хянадаг
\item Хиймэл оюуны үйлчилгээ тасарсан үед нөөц хураангуй болон үндсэн дадлагын даалгавар ашиглах нөөц шийдэлтэй
\item Docker орчин, орчны хувьсагч, аюулгүй байдлын тохиргоотойгоор хөгжүүлэлтээс нэвтрүүлэлт рүү шилжих суурь шийдсэн
\end{enumerate}

\subsection*{Хамрах хүрээ}
Энэхүү ажил нь веб дээр ажиллах цахим сургалтын платформын судалгаа, шинжилгээ, зохиомж, хөгжүүлэлт, туршилт, дүгнэлтийг хамарна. Системийн хэрэглэгч нь үндсэндээ суралцагч бөгөөд програмчлалын агуулга, курс, хичээл, суралцах замнал, түвшин тогтоох шалгалт, түвшин ахих шалгалт, ахицын хураангуй, хиймэл оюуны сургалтын туслах, хамтын орчин, тоглоомын хэсэг гэсэн бүрэлдэхүүнүүдийг багтаана. Туршилт, үр дүнгийн хэсэгт хэрэглэгчийн интерфейсийн бүх дүрс, диаграм, дэлгэцийн зургуудыг тусгайлан оруулах боломжтой байдлаар текстэн тайлбар, хүснэгт, дүгнэлтүүдийг бэлтгэсэн.

\subsection*{Ажлын бүтэц}
Тайлан нь удиртгал, дөрвөн бүлэг, ерөнхий дүгнэлт, хавсралт, ашигласан материалын жагсаалтаас бүрдэнэ. Нэгдүгээр бүлэгт онолын судалгаа, хоёрдугаар бүлэгт системийн шинжилгээ, гуравдугаар бүлэгт зохиомж ба хөгжүүлэлт, дөрөвдүгээр бүлэгт туршилт, үр дүн, хэлэлцүүлгийг авч үзнэ. Хавсралтад системийн чухал кодын хэсэг, API бүтэц, туршилтын жишээ, тайланд оруулах нэмэлт материалын загварыг оруулсан.
